\section{Confluence and Quiescence}
\label{sec:confluence}

The ActiveGraph architecture describes behaviors reacting until quiescence
without a privileged orchestrator.  The working graph is then described as a
deterministic projection of the log.  Projection of a \emph{fixed} log is
deterministic, but the reactive loop can produce different logs under different
legal schedules.  We separate termination from uniqueness of the terminal
projection.

\begin{definition}[Termination at a bound]
For behavior set $B$, initial configuration $C_0$, scheduler $\sigma$, and
bound $m$, $\settle_m(B,C_0,\sigma)$ terminates if it reaches Quiescent or
BlockedAwaitingEffect in at most $m$ steps.  StepBoundReached is evidence of
bounded nontermination, not a proof of infinite reduction.
\end{definition}

\begin{definition}[Schedule confluence under an observation]
A behavior set is confluent for $C_0$ and observation $P$ when every legal
scheduler execution that reaches a quiescent configuration reaches a unique
normal form modulo $\simeq_P$.  Formally, if
$C_0\rightarrow^*_{\sigma_1} C_1$ and
$C_0\rightarrow^*_{\sigma_2} C_2$, and both $C_1$ and $C_2$ are quiescent,
then $C_1\simeq_P C_2$.
\end{definition}

This is Church--Rosser-shaped, but the current artifact does not prove a
Church--Rosser theorem.  It searches scheduler policies, generates restricted
families, and preserves witnesses when the proposition fails.  Newman's lemma
connects termination plus local confluence to confluence in an abstract
rewriting system \cite{newman1942}; our first counterexample removes the
termination premise before the second examines uniqueness of quiescent forms.

\subsection{A1: general termination is false}

\begin{counterexample}[Self-triggering emission]
Let a behavior trigger on signal \code{loop} and emit the same signal.  Seed the
pending queue with \code{loop}.  Every processed signal enables the behavior,
and every firing enqueues another signal.  No quiescent or effect-blocked
configuration is reached.  For every generated bound from 1 through 100, the
artifact reports StepBoundReached at that bound.
\end{counterexample}

The example is intentionally small.  It shows that purity of the reducer and
append-only logging do not imply termination of a reactive runtime.  Detecting
a repeated finite fingerprint would not solve the general problem: generated
identities and sequence counters can make configurations syntactically fresh
while their domain behavior cycles.  The implementation therefore avoids a
heuristic \code{CycleDetected} verdict.

Two natural restoring contracts remain candidates.  One assigns every
internally emitted event a well-founded rank that strictly decreases along
trigger edges.  The other requires an acyclic graph from event kinds to
behavior emissions.  Both would exclude the fixture, but neither is claimed
minimal, and both need refinement for state-dependent triggers.

\subsection{A2: general confluence is false}

\begin{counterexample}[Overlapping writers]
Behaviors $\alpha$ and $\beta$ are both enabled by \code{start}.  Each emits a
write to fact \code{winner}; $\alpha$ writes 1 and $\beta$ writes 2.  Under
canonical activation order the emitted events are processed so that the
terminal winner is 2.  Reversing activation order yields 1.  Both executions
quiesce, but their projected states differ.
\end{counterexample}

FsCheck generalizes this witness over generated integer values by constructing
distinct adjacent values.  Across the generated family, canonical and reverse
activation schedules always disagree.  The write-set diagnostic correctly
reports \code{winner} as a conflict.

It is tempting to restore confluence by requiring declared write sets to be
pairwise disjoint.  That condition works for a deliberately restricted family:
single-trigger behaviors that each unconditionally write a distinct key and do
not read other behaviors' regions or emit cross-triggering signals.  The suite
generates up to eight such writers and checks all four scheduler policies; their
projected states agree.  The structural reason is stronger than those four
samples: every generated key has exactly one unconditional writer, and no
writer reads or triggers another, so any permutation applies the same set of
commuting updates to distinct map entries.  The four policies are executable
regression coverage for that restricted family, not an enumeration of arbitrary
runtime interleavings.

Declared disjointness is not sufficient in general.

\begin{counterexample}[Read/trigger interference with disjoint writes]
Two behaviors, \code{left} and \code{right}, trigger on \code{start} and emit
signals \code{left-ready} and \code{right-ready}.  Their declared write sets are
empty.  A third behavior triggers on \code{left-ready}, reads whether
\code{right-ready} is already in state, and writes the only fact
\code{observed-right}.  All declared write sets are disjoint.  Under canonical
event order the observer records 0; under reverse event order it records 1.
\end{counterexample}

The second witness matters because a write-write conflict detector would
certify it.  The actual dependency crosses four surfaces: an emission, an event
selection, trigger eligibility, and a state read.  Any useful general restoring
condition must constrain at least behavior reads, behavior writes, emitted
event kinds, and the trigger relation.  Under-declared dependencies must either
be rejected or made observable by instrumentation.

\begin{figure}[ht]
\centering
\small
\begin{tabularx}{0.96\textwidth}{@{}p{0.18\textwidth}Y Y@{}}
\toprule
Surface & Canonical event order & Reversed event order \\
\midrule
Emission & \code{left-ready}, then \code{right-ready} & \code{right-ready}, then \code{left-ready} \\
Trigger & Observer becomes eligible before right is projected & Observer becomes eligible after right is projected \\
Read & \code{right-ready} absent & \code{right-ready} present \\
Write & \code{observed-right}=0 & \code{observed-right}=1 \\
\bottomrule
\end{tabularx}
\caption{Counterexample 3 crosses emission, event selection, trigger
eligibility, and a state read despite disjoint declared write sets.}
\label{fig:read-trigger-interference}
\end{figure}

\subsection{Production ordering masks the question}

The ActiveGraph v1.10.0 scheduler uses a FIFO queue, evaluates matching
behaviors in registration order, and dispatches them sequentially.  Work due at
the same delayed tick remains FIFO.  This supplies a reproducible
\emph{ProductionOrder}, and it may be entirely appropriate as an operational
contract.  It is not evidence that alternate legal orders have the same normal
form.

The distinction is practical.  ActiveGraph computes the match set before the
handler loop, so an earlier handler cannot change which behaviors are already
enabled by that trigger.  It does, however, rebuild each handler's view against
the current graph.  An earlier registered handler can emit an event whose
projection is applied synchronously and changes state read by a later handler
of the same original event; behavior matching for the emitted event occurs at a
later FIFO step.  Relation enumeration adds another ordering seam when a
backend query does not specify a canonical order.

The executable kernel deliberately snapshots state across activations enabled
by one trigger, so it does not reproduce this per-invocation production read
timing.  Its tests vary event and activation order directly, and its
read/trigger witness uses distinct later events.  The production seam is a
source-audited limit on generalization, not a property claimed as validated by
the kernel.  A stable FIFO queue can hide either form of dependence in ordinary
replay tests.

\paragraph{Result.}
The strongest supported conclusion is negative and conditional.  General
termination and general confluence are false.  Isolated disjoint writers
converge for the checked family, but declared writes alone do not restore
confluence.  Deriving and validating a general non-interference contract is
future work; this draft does not present a minimal side condition.
